% ===================== sec_conclusion.tex =====================
% NOTE: the future-work items correspond to the limitations in Section 8. Do not promise a release of data or code here unless the release is decided.
\section{Conclusion and Future Work}\label{sec:conclusion}

\subsection{Conclusion}
We presented MedViHall, a benchmark in which twelve vision-language models, nine general-purpose and three domain-adapted, answer four types of query on images from chest radiography, dermoscopy, histopathology and fundus photography. Hallucination is measured at the level of the claim with a taxonomy of seven classes, a severity-weighted rate (\SMHR{}), the share of responses with a fabrication or omission (\CHR{}) and a reference-normalised index of specialisation (\SRI{}). Uncertainty is estimated by resampling images, and all comparisons were audited against the raw result files.

Within the limits of the design, five general-purpose models formed the group with the lowest \SMHR{}, and no order is supported among the first three. Each domain-adapted model had a higher \SMHR{} than the best general-purpose model in each of the four modalities, by a factor between \PH{1.31} and \PH{1.85}, and this result was stable under the sensitivity analyses carried out. The comparison does not isolate the effect of medical adaptation. The level and the composition of the error depend on the modality. The two scores \SMHR{} and \CHR{} ordered the models differently, retrieval-augmented prompting could lower the first while raising the second, and the only mitigation effect that held against both baselines was the reduction of \SMHR{} for one domain-adapted model under contrastive decoding. Stated confidence was far above accuracy for every model examined.

The conclusions are limited by a defect in the radiology arm, by reliance on one automated judge, by author-assigned class weights, and by the absence of clinician annotation; the audit of the labelling pipeline against simulated annotators shows moderate agreement for fabrication and does not support the class-specific labels. We regard the benchmark as a tool for measuring and comparing hallucination under one protocol, and its results as a description of the models and the protocol, not as a statement of clinical safety.

\subsection{Future work}
The following steps would address the main limitations, in the order of their importance for the conclusions.
\begin{enumerate}
  \item \emph{Repair of the radiology arm.} The radiographs should be re-extracted with a match on the exact path, the identity of the images should be verified by checksum, and the radiology queries should be run again for all models. This would restore the radiology arm and allow the transfer of rankings to and from radiology to be assessed.
  \item \emph{Clinician annotation and weights.} A sample of responses should be annotated by clinicians of the corresponding specialities, with several readers per item and a stated protocol, in order to validate the claim-level labels, and the class weights should be elicited from clinicians. The sensitivity of the conclusions to elicited weights can then replace the present test of perturbations.
  \item \emph{Independent judges.} The labelling should be repeated with a judge from another model family and with a larger judge, and the disagreement should be analysed by evaluated model, in order to test for a preference of the judge for similar text.
  \item \emph{Base--specialist pairs.} The models should include pairs of a general-purpose model and the model adapted from it, with full-precision inference, so that the effect of medical adaptation can be separated from differences in size and architecture.
  \item \emph{Repeated sampling and prompt variation.} Several responses per query and several wordings of each query would give a variance that is attributable to the sampling and to the prompt, and would settle the order within the leading group where it can be settled.
  \item \emph{Mitigation.} The strategies should be compared with a baseline generated under identical decoding, with the \CHR{} reported as a primary outcome next to the \SMHR{}, with a knowledge base that clinicians have reviewed, and with the fine-tuning arm that was planned. The localisation of the stage at which a hallucination arises, which is the subject of ClinHallu \citep{yang2026clinhallu}, could be combined with the present taxonomy.
  \item \emph{Wider coverage.} The benchmark could be extended to computed tomography, magnetic resonance imaging and ultrasound, to several images per case, to other languages, and to the generation of complete reports, using the same claim-level protocol.
\end{enumerate}
